% !TEX program = xelatex
\documentclass[12pt,a4paper]{article}
\usepackage[T1]{fontenc}
% --- XeLatex/LuaLatex ---
\usepackage[no-math]{fontspec}
\IfFontExistsTF{Fira Sans}{
	\setmainfont{Fira Sans}
}
\usepackage[bahasa]{babel}
\usepackage{booktabs}
%\renewcommand{\familydefault}{\sfdefault}
\usepackage{geometry}
\usepackage{tabularx}
\usepackage{graphicx}
\usepackage{float}
\usepackage{xcolor}
\usepackage{hyperref}
\usepackage{subcaption}
%\usepackage{subfigure}
\hypersetup{
	colorlinks=true,  % Colours links instead of boxes
	urlcolor=blue,    % Colour for external hyperlinks
	linkcolor=blue,   % Colour of internal links (sections, etc.)
	citecolor=red,    % Colour of citations
	bookmarks=true,   % Adds PDF bookmarks
	hypertexnames=true
}
\geometry{left=2cm, right=2cm, top=2cm}
\title{\textbf{Panduan \textit{Setting} Kamera}}
\author{Agus T. P. Jatmiko}
\hyphenation{me-la-ku-kan peng-amat-an deng-an di-re-ko-men-da-si-kan mem-pe-nga-ruh-i ren-tan peng-ambil-an}
\begin{document}
\maketitle
\noindent\rule{\linewidth}{0.4pt}

\setcounter{section}{1}
\subsection*{Catatan Penting}
\begin{itemize}
	\item Setiap kamera memiliki set parameter yang unik. Gunakan dokumen ini sebagai panduan umum dan sesuaikan dengan kamera yang digunakan.
	\item \textcolor{red}{Jangan mengubah parameter dalam tab \textbf{Image Controls} karena akan mempengaruhi kalibrasi.} Setting default adalah: Gamma=1, Contrast=0, Brightness=0, Timestamp Frames=Off.
	\item Pada tahap ini \textbf{Preprocessing} tidak digunakan. Koreksi FLAT (jika ada) akan dilakukan terpisah.
\end{itemize}

\subsection*{1. Capture Format and Area}
\begin{center}
	\includegraphics[width=0.4\linewidth]{capture_format_area}
\end{center}

\begin{table}[H]
	\centering
	\begin{tabular}{l l p{8cm}}
		Colour Space & : & \textbf{MONO16}\\
		Capture Area & : & \textbf{640 x 480} atau \textbf{320 x 240} (cukup kecil namun masih ada setidaknya satu atau beberapa bintang lain dalam frame)\\
		Binning & : & \textbf{1 x 1}\\
		Output Format & : & \textbf{FITS files (*.fits)} (\textcolor{red}{Penting! direkomendasikan}) atau \textbf{ADV files (*.adv)}\\
	\end{tabular}
\end{table}

\subsection*{2. Camera Controls}
\subsubsection*{2.1. Kamera non--GPS (contoh: ZWO ASI120MM)} 
\begin{center}
	\includegraphics[width=0.4\linewidth]{"camera controls"}
\end{center}

\begin{table}[H]
	\centering
	\begin{tabular}{l l p{8cm}}
		Exposure & : & \textbf{15-50 ms (LX Mode \underline{tidak} dicentang)}\\
		&& Catatan: lakukan tes dalam rentang nilai ini untuk mendapatkan hasil terbaik.\\
		Gain & : & \textbf{70-75 (sesuaikan untuk setiap kamera, sekitar 70-75\% dari total gain)} \\
		Frame Rate Limit & : & \textbf{Maximum}\\
		Turbo USB & : & \textbf{100}\\
		Overclock & : & \textbf{0}\\
		High Speed Mode & : & \textbf{On}\\
	\end{tabular}
\end{table}

Parameter lain biarkan dalam keadaan default.

\subsubsection*{2.2. Kamera GPS (contoh: QHY 174M GPS)}
\begin{center}
	\includegraphics[width=0.4\linewidth]{"camera controls_gps"}
\end{center}

\begin{table}[H]
	\centering
	\begin{tabular}{l l p{8cm}}
		Exposure & : & \textbf{15-50 ms (LX Mode \underline{tidak} dicentang)}\\
		&& Catatan: lakukan tes dalam rentang nilai ini untuk mendapatkan hasil terbaik.\\
		Gain & : & \textbf{300 (sekitar 70-75\% dari total gain)} \\
		Frame Rate Limit & : & \textbf{Maximum}\\
		Amp Noise Reduction & : & \textbf{Auto}\\
		Offset & : & \textbf{100}\\
		USB Traffic & : & \textbf{40}\\
		&& Catatan: semakin tinggi nilai USB Traffic, transfer data semakin cepat, namun rentan terhadap \textit{drop frames}. Jika ada kendala koneksi atau \textit{drop frames}, kurangi nilai ini sampai hasilnya stabil.\\
		Enable Live Broadcast & : & \textbf{Off}\\
		Force Still Mode & : & \textbf{Off}\\
	\end{tabular}
\end{table}

\subsection*{3. Image Controls}
\begin{figure}[H]
	\centering
	\begin{minipage}[t]{0.45\textwidth}
		\includegraphics[width=\linewidth]{"image controls"}
		\caption*{ZWO ASI120MM}
	\end{minipage}\hfil
	\begin{minipage}[t]{0.45\textwidth}
		\includegraphics[width=\linewidth]{"image controls gps"}
		\caption*{QHY 174M GPS}
	\end{minipage}
\end{figure}

Semua \textit{setting} dalam keadaan default seperti pada gambar.
%\begin{table}[H]
%	\centering
%	\begin{tabular}{l l p{8cm}}
%		Timestamp frames & : & \textbf{Off}\\
%	\end{tabular}
%\end{table}

\subsection*{4. GPS Controls (untuk QHY 174M GPS)}
\begin{center}
	\includegraphics[width=0.4\linewidth]{"gps control"}
\end{center}
\begin{table}[H]
	\centering
	\begin{tabular}{l l p{8cm}}
			GPS & : & \textbf{On}\\
			Show Data & : & \textbf{On} (digunakan untuk cek status GPS dan sinkronisasi jam komputer sebelum pengambilan data. )\\
		\end{tabular}
\end{table}
Jangan ubah parameter lainnya. Biarkan dalam keadaan default.

Jika GPS sudah dalam keadaan \textit{locked} (tekan tombol \textbf{Show Data} jika jendela status GPS belum terlihat), tekan tombol \textbf{Set PC Time to GPS Time} sesaat sebelum melakukan pengambilan data.
\begin{center}
	\includegraphics[width=0.4\linewidth]{"gps status"}
\end{center}

\subsection*{5. Thermal Controls}
\begin{center}
	\includegraphics[width=0.4\linewidth]{"thermal controls"}
\end{center}
\begin{table}[H]
	\centering
	\begin{tabular}{l l p{8cm}}
		Cooler Power & : & \textbf{Auto}\\
		Target Temperature & : & \textbf{0}\\
	\end{tabular}
\end{table}
\end{document}